% =====================================================================
%  20.3  Range Assign (Paint)                          sec:c15-assign
%  Exports: prop:c15-override
%           lst:c15-assign  lst:c15-915-soln  prob:c15-pe
%  The assign tag on its own -- the update Ch.18's difference BIT cannot
%  express because it has no inverse.  Aggregate under assignment:
%    sum = value * node-length;  uniform node <=> every block shares a value.
%  CF 915E (live-verified on codeforces.com): n days all working, q orders
%    (l,r,k): k=1 -> days l..r become non-working, k=2 -> working; print the
%    number of working days after EACH order.  n<=1e9 (coordinate compress),
%    q<=3e5.  Sample: n=4 q=6, orders (1,2,1)(3,4,1)(2,3,2)(1,3,2)(2,4,1)(1,4,2)
%    -> 2 0 2 3 1 4  (hand-replayed).
% =====================================================================
\section{Range Assign (Paint)}\label{sec:c15-assign}

\begin{previouslybox}[The update that made the Fenwick tree give up]
Section~\(\secref{c15-compose}\) used \emph{assign} as one tag among several, in service of a MEX problem.  This section promotes range \textbf{assignment} to the headline operation, because it is the single update that Chapter~18 explicitly could not do: to add a value you can subtract it back, but to \emph{set} a position whose old value you do not know there is no inverse to reach for -- the difference trick has nothing to telescope.  The lazy tree handles it with one tag field, and the whole override logic was already paid for in \(\thmref{c15-compose}\).
\end{previouslybox}

\begin{topicintro}
A \term{range assign} (or \term{paint}) update sets every element of \([l,r]\) to one value \(v\).  Under a \emph{sum} aggregate the node fold is a closed form \(\;t_v \leftarrow v\cdot(r_v-l_v+1)\;\) -- you do not need to know the old values, only the length.  The pending tag is simply ``the value I owe my children,'' \(\mathit{tag}_v\in\{v\text{ or }\texttt{none}\}\); a later assignment overrides an earlier one, and \emph{any} assignment overrides an earlier add/flip (the rule proved last section).  The two ideas that make paint-specific problems fast are (i) the \(\mathit{sum}=\mathit{value}\cdot\text{len}\) fold and (ii) \term{uniformity detection}: a node whose whole interval holds one value needs no descent, so a count/max/min query either reads \(t_v\) (uniform) or descends (mixed).

\textbf{Prerequisites.} \(\secref{c15-core}\)'s skeleton and \(\secref{c15-compose}\)'s override law.  The running array \(A\) (here only as a paint demo).  Coordinate compression from \(\secref{c15-compose}\)'s Problem Lab, reused for a \(10^9\)-wide domain.

\textbf{After this section you can}
\begin{itemize}[nosep]
  \item write the assign tag's \code{apply}/\code{push} with the \(\mathit{sum}=v\cdot\text{len}\) fold; \hfill\lv{3}
  \item recognise a \emph{uniform} node and use it to answer ``count of colour \(c\) / are all equal'' queries; \hfill\lv{3}--\lv{4}
  \item handle the mixed \emph{assign}-then-\emph{add} override without a separate case analysis; \hfill\lv{4}
  \item solve a huge-domain paint-and-count input (CF~915E) via compression. \hfill\lv{4}
\end{itemize}
\end{topicintro}

\begin{recognitionbox}[Pattern card -- ``\emph{set} a range to a value, then count / sum / inspect it'']
\reckey{``repaint \([l,r]\) colour \(c\)''; ``make days \(l..r\) free / working''; ``assign \(a_i=v\) for \(i\in[l,r]\)'' followed by ``how many equal \(v\)? / what is the sum? / list the distinct colours?''; a domain too large to store (\(n\) up to \(10^{9}\)) but only \(q\) endpoints matter -- compression; crucially the write is a \emph{set}, not an add.}
\recuse{
  \item paint + count/sum -- the assign tag and \(\mathit{sum}=v\cdot\text{len}\) fold (CF~915E here);
  \item paint + \emph{distinct colours in a range/subtree} -- store a colour \emph{set} (bitmask) per node, OR-merge, assign sets the mask; that is CF~620E in \(\secref{c15-trees}\);
  \item paint \emph{and} add can both occur -- the override decides who wins (assign beats add, \(\thmref{c15-compose}\)).
}
\recnotuse{
  \item the write is an \emph{add}, never a set, and the read is a sum -- Ch.~18's difference BIT;
  \item the write sets a \emph{single} position -- Ch.~19's point update;
  \item the write sets a range to \(\min(a_i,v)\) (a \(\mathit{chmin}\)) -- that is not an assignment to a known value, it is beats, Ch.~23 (\(\secref{c15-strains}\)).
}
One-line test: \emph{you are \textbf{overwriting} a range with a value you know (not incrementing), and later you must read an aggregate of it -- that is the assign tag.}
\end{recognitionbox}

% ---------------------------------------------------------------
\subsection{The assign tag and uniformity}
% ---------------------------------------------------------------

\begin{proposition}{Assignment folds and overrides}{c15-override}
For a range-assign tag on a sum tree: \emph{(fold)} setting a covered node \(v\) of interval length \(\mathit{len}=r_v-l_v+1\) to value \(c\) gives \(t_v=c\cdot\mathit{len}\), independent of the previous values.  \emph{(override)} a new assign \(c\) composed onto any pending tag (assign or add or flip) yields the single tag ``assign \(c\)''.  \emph{(uniformity)} node \(v\) holds one repeated value iff its stored aggregate already determines it uniquely -- for a sum, \(t_v/\mathit{len}\) is that value exactly when \(v\) is uniform; storing an explicit ``all-equal'' flag (set when a cover assigns the node, cleared on any partial descent/pull that finds unequal children) makes uniformity \(O(1)\).
\end{proposition}
\begin{proofbox}
The fold is arithmetic: the sum of \(\mathit{len}\) copies of \(c\) is \(c\cdot\mathit{len}\).  The override is the assign case of \(\thmref{c15-compose}\) -- a set erases history, so whatever the child owed, it now owes only ``set to \(c\).''  For uniformity: after a cover-assign the whole interval is one value, so the flag is set; a later partial update \code{push}es (children inherit the value, staying uniform) then edits one side, and the \code{pull} recomputes \(t_v\) and can only conclude uniform when both children report uniform \emph{with the same value}.  Hence the flag tracks exactly the uniform nodes, in \(O(1)\) per visited node.
\end{proofbox}

\begin{lstlisting}[caption={range-assign tag: the fold is value*len, the merge is a plain overwrite}, label={lst:c15-assign}]
// tag[v]: the value owed to v's children, or -1 (NONE) if clean.
void apply_set(int v, int tl, int tr, long long c){
    t[v]   = c * (tr - tl + 1);   // FOLD: sum = value * length, old values gone
    tag[v] = c;                   // STASH: children now owe "set to c"
    uni[v] = true;                // the whole interval is c -> uniform
}
void push(int v, int tl, int tr){
    if (tag[v] == -1) return;      // NONE: nothing owed
    int tm = (tl + tr) / 2;
    apply_set(2*v,   tl,   tm, tag[v]);
    apply_set(2*v+1, tm+1, tr, tag[v]);
    tag[v] = -1;                   // debt passed down
}
// rupdate / rquery are the SAME as lst:c15-rupdate / lst:c15-rquery, only
// replace "apply(v,tl,tr,d)" by "apply_set(v,tl,tr,c)" and drop the ADD merge.
// Uniformity check in a "count value c" query:
//   if (uni[v] && value_of(v)==c) return (tr-tl+1);   // whole block is c
\end{lstlisting}

\noindent\textbf{Assign \emph{and} add in one tree.}  When both operations occur (paint, then increment some of it), the pending tag becomes a small sum type: ``last assign \(c\) (or none)'' \emph{plus} ``pending add \(d\) since then.''  \code{apply\_set} clears the pending add (override), \code{apply\_add} bumps it only if no assign is waiting on top, and the fold is \(t_v \leftarrow c\cdot\text{len}+d\) when an assign is the base.  This is exactly the affine/override composition of \(\secref{c15-compose}\); nothing new in the traversal.

\begin{insightbox}[Why paint is the poster child for lazy, not BIT]
A difference array records \(a_i-a_{i-1}\).  To set \(a_i:=c\) you would have to know the \emph{old} \(a_i\) to write the new difference -- and it is buried under a thousand earlier updates.  The lazy tag sidesteps the past entirely: it stores the \emph{destination} \(c\), and the fold \(c\cdot\text{len}\) never asks where the values came from.  That asymmetry -- known destination vs known increment -- is precisely the invertibility line Ch.~18 drew and Ch.~20 erases.
\end{insightbox}

% ---------------------------------------------------------------
\subsection{Problem Lab: CF 915E Physical Education Lessons}
% ---------------------------------------------------------------

\solvedproblem{c15-pe}{Physical Education Lessons}{CF 915E}
  {https://codeforces.com/problemset/problem/915/E}{CF 2300}
  {range assign over compressed days; aggregate = count of working days; override-only tag}

\begin{problemstatement}[CF 915E Physical Education Lessons]
There are \(n\) days (numbered \(1..n\)), initially all \emph{working}.  Staff publish \(q\) orders; order \((l,r,k)\) means: if \(k=1\), days \(l..r\) become \emph{non-working}; if \(k=2\), days \(l..r\) become \emph{working}.  Each order fully overrides whatever those days were before.  After \emph{each} order print the number of working days.

\textbf{Input.} First line \(n\) (\(1\le n\le 10^{9}\)); second line \(q\) (\(1\le q\le 3\cdot10^{5}\)); then \(q\) lines \(l\ r\ k\) with \(1\le l\le r\le n\), \(k\in\{1,2\}\).

\textbf{Output.} After each order, one line: the working-day count.

\textbf{Example (official).}
\begin{verbatim}
Input:
4
6
1 2 1
3 4 1
2 3 2
1 3 2
2 4 1
1 4 2

Output:
2
0
2
3
1
4
\end{verbatim}
\end{problemstatement}

\begin{attemptgate}[Attempt it yourself first -- \lv{4}, give yourself 30 minutes]
\(n\) up to \(10^9\): you cannot allocate a per-day array.  And an order is a \emph{set} (working or free), not a toggle -- a BIT cannot express it.
\peek{how many distinct \emph{boundaries} appear across all orders?}
\nudge{only the endpoints \(l\) and \(r+1\) ever split a day-run, so between consecutive sorted endpoints every day shares its fate -- compress those \(\le 2q{+}2\) points into blocks, exactly as in \(\secref{c15-compose}\)'s MEX.}
\fullhint{Over the compressed blocks build an assign tree whose aggregate is ``working days in this node'' and whose tag is \(\{-1,0,1\}\) (free / working / none).  Set \(\mathit{val}=(k{=}2?1:0)\); a cover assigns \(t_v=\mathit{val}\cdot\mathit{len}_v\).  Map order \([l,r]\) to blocks \([\mathrm{idx}(l),\mathrm{idx}(r{+}1){-}1]\).  Initially everything is working: \(t_{\text{root}}=n\).  Print \(t[1]\) after each order -- the whole answer is the root sum, no query walk needed.}
\end{attemptgate}

\solstep{Understanding}
Every order overwrites a day-range with ``working'' or ``free'' -- a pure range-assign.  The quantity demanded is a single global number, ``how many days are working,'' which is the sum-tree aggregate at the root (each block contributes its working-day count).  The \(10^9\) range is the only complication; coordinate compression to the \(O(q)\) endpoints makes the tree tiny, and the block-length carries the multiplicities so the sum is exact.

\solstep{Key Observations}
\begin{observation}
The answer after every order is just \(t[1]\) -- the root already stores the total working count.  There is no \emph{range} read at all; this is ``paint + total count,'' the lightest query there is, so the whole cost is the update.
\end{observation}
\begin{observation}
Only endpoints \(l\) and \(r+1\) can change a day's status as a group; between two consecutive sorted endpoints all days receive identical orders.  Compress \(\{1,n{+}1\}\cup\{l_i,r_i{+}1\}\).
\end{observation}
\begin{observation}
The tag is override-only: \(k{=}1\) assigns~0, \(k{=}2\) assigns~1; the fold \(\mathit{sum}=\mathit{val}\cdot\mathit{len}\) uses the block-length, so an interval of \(10^5\) consecutive identical days costs one node.
\end{observation}

\solstep{Think It Through With Me}
\begin{thinkbox}
\deadend{Keep a set of working intervals and edit it}
{A ``keep maximal free/working runs'' (Ch.~6 \texttt{std::set}) approach is the classic non-segtree solution, but each order can split and merge many runs; amortised it works, yet with \(q=3\cdot10^5\) the constant and the erase-loop are fragile.  The compressed assign tree is one clean \(O(\log q)\) per order.}
\deadend{Allocate \(t[1..n]\)}
{\(n=10^9\) \(\Rightarrow\) a \(4n\) tree is \(1.6\cdot10^{10}\) slots -- \mle (and you never read \(99.9\%\) of it).  Compress first.}
\aha{Compress to \(O(q)\) uniform blocks, run a lazy \emph{assign} tree storing working-days per node, and read the root.}
\end{thinkbox}

\solstep{The Approach Ladder}
\begin{center}\renewcommand{\arraystretch}{1.25}
\begin{tabular}{@{}c p{0.20\linewidth} p{0.28\linewidth} l c@{}}
\toprule
\textbf{\#} & \textbf{Approach} & \textbf{Idea} & \textbf{Time} & \textbf{Verdict}\\
\midrule
1 & \lv{3}\ array over days & \(t[1..n]\), assign per day & \(\bigO{qn}\) & \mle \\
2 & \lv{3}\ run-set (Ch.~6) & split/merge working runs & \(\bigO{q\log q}\) & passes \\
3 & \lv{4}\ \textbf{compressed assign} & \(O(\log q)\) assign per order & \(\bigO{q\log q}\) & \ac \\
4 & \lv{4}\ implicit assign tree & grow nodes, no precompress & \(\bigO{q\log n}\) & \ac \\
\bottomrule
\end{tabular}\end{center}
Rung~2 is legitimate and is the same idea as Ch.~6's \texttt{std::set} interval trick; we take rung~3 to keep the lazy engine identical to the rest of the chapter.

\solstep{Solution}
Collect points \(\{1\}\cup\{n+1\}\cup\{l_i\}\cup\{r_i+1\}\), sort unique into \(xs\); block \(i=[xs_i,xs_{i+1}-1]\) has length \(xs_{i+1}-xs_i\).  Build an assign tree over \([0,|xs|-2]\) with \(\mathit{sum}=\mathit{len}\) (all working) and \(\mathit{tag}=\) none.  For order \((l,r,k)\): \(\mathit{val}=(k{=}2)\), map to blocks \([\mathrm{idx}(l),\mathrm{idx}(r{+}1)-1]\), \code{assign\_rng}.  Print \(t[1]\).

\solstep{Proof of Correctness}
\begin{proofbox}
\textbf{Block uniformity.} Every order boundary is a point of \(xs\); hence a block \([xs_i,xs_{i+1}-1]\) is wholly inside or outside every order and all its days always share status \emph{working}/\emph{free}.  Storing one bit per block, weighted by its length, is exact.
\textbf{Assign semantics.} \code{apply\_set} on a covered node sets \(\mathit{sum}=\mathit{val}\cdot\mathit{len}\) (\(\propref{c15-override}\)) and stashes the value so children, when pushed, inherit it -- matching ``order overrides previous status.''  \code{push}/\code{pull} preserve the sum invariant (\(\thmref{c15-lazy}\)).
\textbf{Root count.} \(t[1]=\sum_{\text{blocks}} \mathit{working}_i\cdot\mathit{len}_i\) = number of working days, since the block lengths partition \([1,n]\).  Mapping \([l,r]\) to \([\mathrm{idx}(l),\mathrm{idx}(r{+}1)-1]\) selects exactly the blocks within \([l,r]\).  \textbf{Bounds.} \(|xs|\le 2q+2\); all indices fit \code{int}, day counts fit \code{long long}.
\end{proofbox}

\solstep{Implementation}
\begin{lstlisting}[caption={CF 915E full solution: compressed lazy assign tree; the answer is the root sum}, label={lst:c15-915-soln}]
#include <bits/stdc++.h>
using namespace std;
const int MX = 600005;                 // <= 2*q + 2 boundary points
long long t[4*MX], ln[4*MX];           // working days; days covered by node
int tagv[4*MX];                        // -1 none, 0 free, 1 working
vector<long long> xs;

void build(int v, int l, int r){
    tagv[v] = -1;
    ln[v] = xs[r+1] - xs[l];           // number of days in blocks l..r
    t[v]  = ln[v];                     // initially every day is working
    if (l != r){ int m=(l+r)/2; build(2*v,l,m); build(2*v+1,m+1,r); }
}
void apply_set(int v, int val){
    t[v] = val ? ln[v] : 0;            // fold: sum = val * length
    tagv[v] = val;                     // override the children's owed value
}
void push(int v){
    if (tagv[v] == -1) return;
    apply_set(2*v, tagv[v]); apply_set(2*v+1, tagv[v]);
    tagv[v] = -1;
}
void assign_rng(int v, int l, int r, int ql, int qr, int val){
    if (ql > r || qr < l) return;
    if (ql <= l && r <= qr){ apply_set(v, val); return; }   // cover: fold + stop
    push(v); int m = (l+r)/2;
    assign_rng(2*v,l,m,ql,qr,val); assign_rng(2*v+1,m+1,r,ql,qr,val);
    t[v] = t[2*v] + t[2*v+1];                               // pull
}
int main(){
    long long n; int q;
    scanf("%lld%d", &n, &q);
    vector<array<long long,3>> ord(q);
    xs.push_back(1); xs.push_back(n+1);
    for (int i = 0; i < q; i++){
        long long l,r,k; scanf("%lld%lld%lld", &l,&r,&k);
        ord[i] = {l,r,k}; xs.push_back(l); xs.push_back(r+1);
    }
    sort(xs.begin(), xs.end()); xs.erase(unique(xs.begin(), xs.end()), xs.begin());
    int nseg = (int)xs.size() - 1;
    build(1, 0, nseg-1);
    for (int i = 0; i < q; i++){
        long long l=ord[i][0], r=ord[i][1], k=ord[i][2];
        int L = lower_bound(xs.begin(), xs.end(), l)   - xs.begin();
        int R = lower_bound(xs.begin(), xs.end(), r+1) - xs.begin() - 1;
        assign_rng(1, 0, nseg-1, L, R, k==2 ? 1 : 0);   // k=1 free, k=2 working
        printf("%lld\n", t[1]);                         // root = working days
    }
}
\end{lstlisting}

\solstep{Trace on the Official Sample}
\begin{tracebox}[CF 915E: \(n=4\); orders \((1,2,1),(3,4,1),(2,3,2),(1,3,2),(2,4,1),(1,4,2)\)]
\(xs=(1,2,3,4,5)\): four length-1 blocks \([1],[2],[3],[4]\), start working \(([1,1,1,1])\), root \(=4\).
\begin{center}\renewcommand{\arraystretch}{1.15}
\begin{tabular}{@{}c l l l c@{}}
\toprule
\textbf{order} & \textbf{action} & \textbf{blocks set} & \textbf{days} & \textbf{root}\\
\midrule
\(1\ 2\ 1\) & assign 0 & blocks \(0..1\) & \((0,0,1,1)\) & \(2\)\\
\(3\ 4\ 1\) & assign 0 & blocks \(2..3\) & \((0,0,0,0)\) & \(0\)\\
\(2\ 3\ 2\) & assign 1 & blocks \(1..2\) & \((0,1,1,0)\) & \(2\)\\
\(1\ 3\ 2\) & assign 1 & blocks \(0..2\) & \((1,1,1,0)\) & \(3\)\\
\(2\ 4\ 1\) & assign 0 & blocks \(1..3\) & \((1,0,0,0)\) & \(1\)\\
\(1\ 4\ 2\) & assign 1 & blocks \(0..3\) & \((1,1,1,1)\) & \(4\)\\
\bottomrule
\end{tabular}\end{center}
Printed \(\boxed{2\;0\;2\;3\;1\;4}\) -- the official output.  Each line is the root read in \(O(1)\); only the update walks the tree.
\end{tracebox}

\solstep{Complexity}
\begin{complexitybox}
\(|xs|\le 2q+2\le 6\cdot10^5{+}2\), so \(nseg=O(q)\).  Build \(O(q)\); each order one \(\bigO{\log q}\) assign \(\Rightarrow\) total \(\bigO{q\log q}\approx 1.2\cdot10^7\) node-visits for \(q=3\cdot10^5\).  Memory \(4\cdot nseg\) of \(\{t,\mathit{ln}\}\) \code{long long} + \(\mathit{tagv}\) \code{int} \(\approx 48\) MB.  Time limit is tight (1 s) -- the \(O(1)\) root read matters: no per-order \code{rquery}.
\end{complexitybox}

\solstep{Verification and Edge Cases}
\begin{center}\renewcommand{\arraystretch}{1.2}
\begin{tabular}{@{}p{0.22\linewidth} p{0.30\linewidth} p{0.36\linewidth}@{}}
\toprule
\textbf{Case} & \textbf{Why tricky} & \textbf{Code behaviour}\\
\midrule
\(n=10^{9}\), 1 order & huge domain & \(xs=\{1,l,r{+}1,n{+}1\}\), a few blocks; lengths carry counts \\
re-assign same range & idempotent & override sets the same value, root unchanged \\
overlapping orders & later must win & \code{apply\_set} overwrites the stale tag (\(\propref{c15-override}\)) \\
\(k{=}1\) on all-free & no change & assigns 0 to blocks already 0, sum stays \\
root read & per-order cost & printed straight from \(t[1]\), \(O(1)\) \\
\bottomrule
\end{tabular}\end{center}

\begin{hardeningbox}
\begin{itemize}[nosep,leftmargin=1.4em]
  \item \textbf{Compress \(r+1\), and bound with \(1\) and \(n+1\).}  Missing a boundary merges two days that could be edited apart; missing \(n+1\) drops the trailing always-working run from the count.
  \item \textbf{Fold with the node's \emph{day-length}, not block count.}  \(\mathit{len}_v\) is a \emph{number of days}, so a length-\(10^5\) block contributes \(10^5\) working days, not \(1\).
  \item \textbf{Read \(t[1]\), do not query.}  The asked quantity is global; a \code{rquery} per order adds a needless \(\log\).  (When the task instead asks a \emph{range} count, use \(\coderef{c15-rquery}\).)
  \item \textbf{Assign tag has no ``add'' merge.}  If you later mix in increments, the tag must become the assign-then-add sum type -- do not reuse the pure-add \(\mathit{tag}\mathrel{+}=\) from \(\secref{c15-core}\).
\end{itemize}
\end{hardeningbox}

% ---------------------------------------------------------------
\subsection{Common Pitfalls}
% ---------------------------------------------------------------

\begin{warningbox}[Trap 1: the ``set-to-0 means never-touched'' confusion]
A pending tag of ``assign \(0\)'' is \emph{real work} -- it must reach the children.  Using \(\mathit{tag}=0\) as the ``clean/none'' sentinel (as the add tree does) makes an assign-to-0 silently dropped.  Use a distinct NONE sentinel (\(-1\) above; or a separate boolean) whenever \(0\) is a legal assigned value.
\end{warningbox}

\begin{warningbox}[Trap 2: forgetting a \emph{later} assignment overrides an \emph{earlier} one]
Two orders hitting the same dirty node: the second must \emph{replace} the first's tag, not add to it.  Leaving the old value composed (e.g.\ treating assign as an add) yields counts one order stale -- passing samples where no range is painted twice.
\end{warningbox}

\begin{warningbox}[Trap 3: forgetting to \code{push} before a partial paint]
Same rule as \(\secref{c15-core}\): descending through a node that carries an assign requires pushing it, or the children keep values the parent already overwrote -- the partial paint is applied to a stale base.
\end{warningbox}

\begin{warningbox}[Trap 4: uniformity assumed, never verified]
Answering ``are all of \([l,r]\) colour \(c\)'' by checking only \(t_v/\mathit{len}=c\) is wrong when the node is \emph{mixed} yet happens to average to \(c\) (e.g.\ one \(0\) and one \(2\) average to \(1\)).  Trust the explicit \(\mathit{uni}\) flag (\(\propref{c15-override}\)), not the ratio, unless you keep min and max too and check \(\min=\max\).
\end{warningbox}

% ---------------------------------------------------------------
\subsection{Where to Drill Next}
% ---------------------------------------------------------------

\begin{stuckbox}
\begin{itemize}[nosep,leftmargin=1.4em]
  \item Is your ``none'' tag distinguishable from a legal assigned value (especially \(0\))?
  \item Does a later assign \emph{overwrite} (not combine with) the pending tag?
  \item Is your fold \(c\cdot\text{node-length}\) using the \emph{days}, not the block count?
  \item When you descend past an assigned node, did you \code{push} first?
  \item Did you bound the compression with \(1\) and \(n+1\) so no trailing run is lost?
\end{itemize}
\end{stuckbox}

The pure-array paint here is the same tag CF~444C (\emph{DZY Loves Colors}, in the chapter-end mixed set) drives over an \(O(n)\) domain, and the ``distinct colours in a range/subtree'' upgrade -- store a colour \emph{bitmask} per node, OR-merge, assign sets a single-bit mask -- is exactly CF~620E, which \(\secref{c15-trees}\) solves on a flattened tree.  To self-test uniformity detection: paint a small array with alternating assigns and query ``is the whole range one colour?'' both with the \(\mathit{uni}\) flag and (wrongly) with the sum-ratio, and find the mixed-but-average-equal counterexample.

\paragraph{Back-cites and forwards.}
\(\secref{c15-compose}\) owns the override law and the assign/flip tag; \(\secref{c15-core}\) owns \code{push}-before-descend.  Chapter~18 (\chref{c14}) is the contrast this section finally makes concrete (assignment is the non-invertible update it handed off).  Forward: \(\secref{c15-trees}\) paints \emph{subtrees} (CF~620E) with this tag; \(\secref{c15-strains}\) notes that range \(\mathit{chmin}\) is a ``set to \(\min\),'' \emph{not} a set-to-a-known-value, and so is beats, Ch.~23 (\chref{c29}).

% ---------------------------------------------------------------
\begin{recallbox}[Section 20.3 in 10 seconds]
Range \textbf{assign} (paint) is the update with no inverse -- beyond Ch.~18's difference BIT, home turf of lazy.  The tag is the owed value (with a NONE sentinel distinct from any legal value); \code{apply\_set} folds \(t_v=v\cdot\text{len}\) and overrides any older tag (\(\propref{c15-override}\)); a later paint beats an earlier one, and any paint beats an add/flip.  Track a \emph{uniform} flag for count/are-all-equal queries.  A too-wide domain (\(10^9\)) is compressed to \(O(q)\) endpoints.  Worked: CF~915E (\(\probref{c15-pe}\), \(\coderef{c15-915-soln}\)) -- compressed assign tree, answer \(=\) root sum.
\selftest{
  \item Why can a difference BIT not do ``set \([l,r]\) to \(c\)'' but can do ``add \(d\) to \([l,r]\)''?
  \item Write the fold for sum, min, and ``count of colour \(c\)'' under an assign of value \(v\).
  \item What goes wrong if \(\mathit{tag}=0\) doubles as both ``assign free'' and ``clean''?
  \item In CF~915E, why is the per-order read \(O(1)\) rather than \(O(\log n)\)?
  \item How would you extend the node record to answer ``distinct colours in the range'' (foreshadow CF~620E)?
}
\end{recallbox}
